\documentclass[12pt,a4paper]{book}
\usepackage[utf8]{inputenc}
\usepackage{amsmath,amsfonts,amssymb}
\usepackage{listings}
\usepackage{xcolor}
\usepackage{geometry}
\usepackage{hyperref}
\usepackage{fancyhdr}
\usepackage{graphicx}
\usepackage{titlesec}

\geometry{margin=1in}
\pagestyle{fancy}
\fancyhf{}
\rhead{Python Data Structures - Complete Guide}
\lfoot{Python for Data Science}
\rfoot{\thepage}

% Python code styling
\lstset{
    language=Python,
    basicstyle=\ttfamily\small,
    keywordstyle=\color{blue}\bfseries,
    commentstyle=\color{green!60!black},
    stringstyle=\color{red},
    numbers=left,
    numberstyle=\tiny\color{gray},
    stepnumber=1,
    numbersep=8pt,
    showstringspaces=false,
    breaklines=true,
    frame=single,
    backgroundcolor=\color{gray!10}
}

% Chapter styling
\titleformat{\chapter}[display]
{\normalfont\huge\bfseries}{\chaptertitlename\ \thechapter}{20pt}{\Huge}

\title{\textbf{Python Data Structures\\Complete Guide}}
\author{Python for Data Science - Beginner's Level\\Taught by Mohit Sir}
\date{\today}

\begin{document}

\maketitle

\tableofcontents
\newpage

\chapter*{Preface}
\addcontentsline{toc}{chapter}{Preface}

This comprehensive guide covers all major Python data structures essential for data science and programming. Each chapter provides detailed explanations, practical examples, and best practices.

\textbf{What You'll Learn:}
\begin{itemize}
    \item Core Python data structures: Lists, Tuples, Dictionaries, Sets, Strings
    \item Scientific computing with NumPy arrays
    \item Data analysis and manipulation with Pandas
    \item Performance considerations and optimization techniques
    \item Real-world applications and use cases
\end{itemize}

\textbf{Prerequisites:}
\begin{itemize}
    \item Basic Python syntax knowledge
    \item Understanding of programming concepts
    \item Python 3.x installed on your system
\end{itemize}

\chapter{Introduction to Python Data Structures}

Python provides several built-in data structures that are essential for effective programming and data science. Understanding when and how to use each structure is crucial for writing efficient code.

\section{Overview of Data Structures}

\begin{table}[h]
\centering
\begin{tabular}{|l|c|c|c|l|}
\hline
\textbf{Structure} & \textbf{Ordered} & \textbf{Mutable} & \textbf{Duplicates} & \textbf{Use Case} \\
\hline
List & \checkmark & \checkmark & \checkmark & Dynamic collections \\
Tuple & \checkmark & $\times$ & \checkmark & Immutable sequences \\
Dictionary & \checkmark* & \checkmark & Keys: $\times$ & Key-value mapping \\
Set & $\times$ & \checkmark & $\times$ & Unique collections \\
String & \checkmark & $\times$ & \checkmark & Text processing \\
NumPy Array & \checkmark & \checkmark & \checkmark & Numerical computing \\
Pandas DataFrame & \checkmark & \checkmark & \checkmark & Data analysis \\
\hline
\end{tabular}
\caption{Python Data Structures Comparison}
\end{table}

*Ordered since Python 3.7+

\section{Performance Considerations}

Different data structures have different performance characteristics:

\begin{itemize}
    \item \textbf{Lists}: O(1) append, O(n) search, O(n) insert at beginning
    \item \textbf{Tuples}: O(1) access, O(n) search, immutable (no modifications)
    \item \textbf{Dictionaries}: O(1) average access/insert/delete, O(n) worst case
    \item \textbf{Sets}: O(1) average membership testing, O(n) worst case
    \item \textbf{NumPy Arrays}: Vectorized operations, memory efficient
    \item \textbf{Pandas}: Optimized for data analysis workflows
\end{itemize}

\chapter{Lists}
\documentclass[12pt,a4paper]{article}
\usepackage[utf8]{inputenc}
\usepackage{amsmath,amsfonts,amssymb}
\usepackage{listings}
\usepackage{xcolor}
\usepackage{geometry}
\usepackage{hyperref}
\usepackage{fancyhdr}

\geometry{margin=1in}
\pagestyle{fancy}
\fancyhf{}
\rhead{Python Lists - Complete Guide}
\lfoot{Python for Data Science}
\rfoot{\thepage}

% Python code styling
\lstset{
    language=Python,
    basicstyle=\ttfamily\small,
    keywordstyle=\color{blue}\bfseries,
    commentstyle=\color{green!60!black},
    stringstyle=\color{red},
    numbers=left,
    numberstyle=\tiny\color{gray},
    stepnumber=1,
    numbersep=8pt,
    showstringspaces=false,
    breaklines=true,
    frame=single,
    backgroundcolor=\color{gray!10}
}

\title{\textbf{Python Lists - Complete Guide}}
\author{Python for Data Science - Beginner's Level}
\date{\today}

\begin{document}

\maketitle
\tableofcontents
\newpage

\section{Overview}

Lists are one of the most versatile and commonly used data structures in Python. They are \textbf{ordered}, \textbf{mutable} (changeable), and allow \textbf{duplicate elements}. Lists can store elements of different data types.

\section{List Creation}

\subsection{Basic Creation}

\begin{lstlisting}
# Empty list
empty_list = []
empty_list = list()

# List with elements
numbers = [1, 2, 3, 4, 5]
mixed_list = [1, "hello", 3.14, True]
nested_list = [[1, 2], [3, 4], [5, 6]]

# Using list() constructor
from_string = list("hello")  # ['h', 'e', 'l', 'l', 'o']
from_range = list(range(5))  # [0, 1, 2, 3, 4]
\end{lstlisting}

\subsection{List Comprehension}

\begin{lstlisting}
# Basic comprehension
squares = [x**2 for x in range(10)]

# With condition
even_squares = [x**2 for x in range(10) if x % 2 == 0]

# Nested comprehension
matrix = [[i*j for j in range(3)] for i in range(3)]
\end{lstlisting}

\section{Accessing Elements}

\subsection{Indexing}

\begin{lstlisting}
fruits = ["apple", "banana", "cherry", "date"]

# Positive indexing (0-based)
print(fruits[0])    # "apple"
print(fruits[2])    # "cherry"

# Negative indexing (-1 is last element)
print(fruits[-1])   # "date"
print(fruits[-2])   # "cherry"
\end{lstlisting}

\subsection{Slicing}

\begin{lstlisting}
numbers = [0, 1, 2, 3, 4, 5, 6, 7, 8, 9]

# Basic slicing [start:end:step]
print(numbers[2:7])     # [2, 3, 4, 5, 6]
print(numbers[:5])      # [0, 1, 2, 3, 4]
print(numbers[5:])      # [5, 6, 7, 8, 9]
print(numbers[::2])     # [0, 2, 4, 6, 8]
print(numbers[::-1])    # [9, 8, 7, 6, 5, 4, 3, 2, 1, 0]
\end{lstlisting}

\section{Modifying Lists}

\subsection{Adding Elements}

\begin{lstlisting}
fruits = ["apple", "banana"]

# append() - Add single element at end
fruits.append("cherry")

# insert() - Add element at specific position
fruits.insert(1, "orange")

# extend() - Add multiple elements
fruits.extend(["grape", "mango"])

# Using + operator
more_fruits = fruits + ["kiwi", "peach"]
\end{lstlisting}

\subsection{Removing Elements}

\begin{lstlisting}
fruits = ["apple", "banana", "cherry", "banana", "date"]

# remove() - Remove first occurrence
fruits.remove("banana")

# pop() - Remove and return element by index
removed = fruits.pop(1)

# del statement
del fruits[0]

# clear() - Remove all elements
fruits.clear()
\end{lstlisting}

\section{List Methods}

\subsection{Search Methods}

\begin{lstlisting}
fruits = ["apple", "banana", "cherry", "banana", "date"]

# index() - Find first occurrence
idx = fruits.index("banana")

# count() - Count occurrences
count = fruits.count("banana")

# in operator - Check if element exists
print("apple" in fruits)     # True
\end{lstlisting}

\subsection{Sorting Methods}

\begin{lstlisting}
numbers = [3, 1, 4, 1, 5, 9, 2, 6]

# sort() - Sort in place
numbers.sort()

# sorted() - Return new sorted list
sorted_list = sorted(numbers)

# Custom sorting with key function
words = ["banana", "pie", "Washington", "book"]
words.sort(key=len)  # Sort by length
\end{lstlisting}

\section{Performance Considerations}

\begin{table}[h]
\centering
\begin{tabular}{|l|c|l|}
\hline
\textbf{Operation} & \textbf{Time Complexity} & \textbf{Notes} \\
\hline
Access by index & O(1) & Direct access \\
Search & O(n) & Linear search \\
Append & O(1) amortized & May need to resize \\
Insert at beginning & O(n) & Shifts all elements \\
Delete by index & O(n) & May shift elements \\
Sort & O(n log n) & Timsort algorithm \\
\hline
\end{tabular}
\caption{List Operations Time Complexity}
\end{table}

\section{Common Use Cases}

\subsection{Stack Implementation}

\begin{lstlisting}
stack = []

# Push
stack.append(1)
stack.append(2)
stack.append(3)

# Pop
top = stack.pop()  # Returns 3
\end{lstlisting}

\subsection{Matrix Operations}

\begin{lstlisting}
# 2D matrix using nested lists
matrix = [
    [1, 2, 3],
    [4, 5, 6],
    [7, 8, 9]
]

# Access element
print(matrix[1][2])  # 6

# Transpose matrix
transposed = [[row[i] for row in matrix] for i in range(len(matrix[0]))]
\end{lstlisting}

\section{Best Practices}

\begin{itemize}
    \item Use list comprehensions for simple transformations
    \item Avoid modifying lists while iterating over them
    \item Use appropriate data types for better performance
    \item Consider using \texttt{collections.deque} for frequent insertions/deletions at both ends
    \item Use \texttt{enumerate()} when you need both index and value
\end{itemize}

\section{Summary}

Lists are fundamental Python data structures that provide:
\begin{itemize}
    \item \textbf{Flexibility}: Store any type of data
    \item \textbf{Mutability}: Can be modified after creation
    \item \textbf{Order}: Maintain insertion order
    \item \textbf{Indexing}: Fast access by position
    \item \textbf{Rich methods}: Extensive built-in functionality
\end{itemize}

Master lists to build a strong foundation for Python programming and data science!

\end{document}

\chapter{Tuples}
\documentclass[12pt,a4paper]{article}
\usepackage[utf8]{inputenc}
\usepackage{amsmath,amsfonts,amssymb}
\usepackage{listings}
\usepackage{xcolor}
\usepackage{geometry}
\usepackage{hyperref}
\usepackage{fancyhdr}

\geometry{margin=1in}
\pagestyle{fancy}
\fancyhf{}
\rhead{Python Tuples - Complete Guide}
\lfoot{Python for Data Science}
\rfoot{\thepage}

\lstset{
    language=Python,
    basicstyle=\ttfamily\small,
    keywordstyle=\color{blue}\bfseries,
    commentstyle=\color{green!60!black},
    stringstyle=\color{red},
    numbers=left,
    numberstyle=\tiny\color{gray},
    stepnumber=1,
    numbersep=8pt,
    showstringspaces=false,
    breaklines=true,
    frame=single,
    backgroundcolor=\color{gray!10}
}

\title{\textbf{Python Tuples - Complete Guide}}
\author{Python for Data Science - Beginner's Level}
\date{\today}

\begin{document}

\maketitle
\tableofcontents
\newpage

\section{Overview}

Tuples are \textbf{ordered}, \textbf{immutable} (unchangeable), and allow \textbf{duplicate elements}. They are similar to lists but cannot be modified after creation, making them ideal for storing data that shouldn't change.

\section{Tuple Creation}

\subsection{Basic Creation}

\begin{lstlisting}
# Empty tuple
empty_tuple = ()
empty_tuple = tuple()

# Tuple with elements
numbers = (1, 2, 3, 4, 5)
mixed_tuple = (1, "hello", 3.14, True)

# Single element tuple (comma is required!)
single_element = (42,)  # Without comma, it's just parentheses
single_element = 42,    # Alternative syntax

# Without parentheses (tuple packing)
coordinates = 10, 20, 30

# Nested tuples
nested_tuple = ((1, 2), (3, 4), (5, 6))
\end{lstlisting}

\subsection{Using tuple() Constructor}

\begin{lstlisting}
# From list
from_list = tuple([1, 2, 3, 4])  # (1, 2, 3, 4)

# From string
from_string = tuple("hello")  # ('h', 'e', 'l', 'l', 'o')

# From range
from_range = tuple(range(5))  # (0, 1, 2, 3, 4)
\end{lstlisting}

\section{Accessing Elements}

\subsection{Indexing and Slicing}

\begin{lstlisting}
fruits = ("apple", "banana", "cherry", "date")

# Positive indexing
print(fruits[0])    # "apple"
print(fruits[2])    # "cherry"

# Negative indexing
print(fruits[-1])   # "date"
print(fruits[-2])   # "cherry"

# Slicing
numbers = (0, 1, 2, 3, 4, 5, 6, 7, 8, 9)
print(numbers[2:7])     # (2, 3, 4, 5, 6)
print(numbers[::-1])    # (9, 8, 7, 6, 5, 4, 3, 2, 1, 0)
\end{lstlisting}

\section{Tuple Methods}

\subsection{Built-in Methods}

\begin{lstlisting}
fruits = ("apple", "banana", "cherry", "banana", "date")

# count() - Count occurrences of an element
banana_count = fruits.count("banana")  # 2

# index() - Find first occurrence of an element
banana_index = fruits.index("banana")  # 1

# index() with start parameter
banana_index_after = fruits.index("banana", 2)  # 3
\end{lstlisting}

\subsection{Membership Testing}

\begin{lstlisting}
fruits = ("apple", "banana", "cherry")

# in operator
print("apple" in fruits)     # True
print("grape" in fruits)     # False

# not in operator
print("grape" not in fruits) # True
\end{lstlisting}

\section{Tuple Operations}

\subsection{Mathematical Operations}

\begin{lstlisting}
# Concatenation
tuple1 = (1, 2, 3)
tuple2 = (4, 5, 6)
combined = tuple1 + tuple2  # (1, 2, 3, 4, 5, 6)

# Repetition
repeated = (1, 2) * 3  # (1, 2, 1, 2, 1, 2)

# Length
print(len((1, 2, 3, 4, 5)))  # 5

# Min, Max, Sum (for numeric tuples)
numbers = (3, 1, 4, 1, 5, 9, 2, 6)
print(min(numbers))  # 1
print(max(numbers))  # 9
print(sum(numbers))  # 31
\end{lstlisting}

\section{Tuple Unpacking}

\subsection{Basic Unpacking}

\begin{lstlisting}
# Assign tuple elements to variables
point = (10, 20)
x, y = point
print(f"x: {x}, y: {y}")  # x: 10, y: 20

# Multiple assignment
person = ("Alice", 25, "Engineer")
name, age, job = person
\end{lstlisting}

\subsection{Advanced Unpacking}

\begin{lstlisting}
# Using * for collecting multiple elements
numbers = (1, 2, 3, 4, 5)
first, *middle, last = numbers
print(f"First: {first}")    # First: 1
print(f"Middle: {middle}")  # Middle: [2, 3, 4]
print(f"Last: {last}")      # Last: 5

# Swapping variables
a, b = 10, 20
a, b = b, a  # Swap values
\end{lstlisting}

\section{Named Tuples}

\begin{lstlisting}
from collections import namedtuple

# Define a named tuple class
Point = namedtuple('Point', ['x', 'y'])
Person = namedtuple('Person', ['name', 'age', 'city'])

# Create instances
p1 = Point(10, 20)
person1 = Person("Alice", 25, "New York")

# Access by name or index
print(p1.x, p1.y)        # 10 20
print(person1.name)      # Alice

# Named tuple methods
point_dict = p1._asdict()  # Convert to dictionary
p2 = p1._replace(x=30)     # Create new with changed field
\end{lstlisting}

\section{Performance Considerations}

\begin{table}[h]
\centering
\begin{tabular}{|l|c|l|}
\hline
\textbf{Operation} & \textbf{Time Complexity} & \textbf{Notes} \\
\hline
Access by index & O(1) & Direct access \\
Search & O(n) & Linear search \\
Count & O(n) & Must check all elements \\
Concatenation & O(n+m) & Creates new tuple \\
Slicing & O(k) & k is slice length \\
\hline
\end{tabular}
\caption{Tuple Operations Time Complexity}
\end{table}

\subsection{Memory Efficiency}

Tuples are more memory efficient than lists:
\begin{itemize}
    \item Smaller memory footprint
    \item Faster creation and iteration
    \item Can be used as dictionary keys (hashable)
\end{itemize}

\section{Common Use Cases}

\subsection{Coordinates and Points}

\begin{lstlisting}
# 2D coordinates
point_2d = (10, 20)
x, y = point_2d

# 3D coordinates
point_3d = (10, 20, 30)
x, y, z = point_3d

# RGB colors
red = (255, 0, 0)
green = (0, 255, 0)
blue = (0, 0, 255)
\end{lstlisting}

\subsection{Function Return Values}

\begin{lstlisting}
def get_name_age():
    """Return multiple values as a tuple"""
    return "Alice", 25

def calculate_stats(numbers):
    """Return statistics as a tuple"""
    return min(numbers), max(numbers), sum(numbers) / len(numbers)

# Usage
name, age = get_name_age()
minimum, maximum, average = calculate_stats([1, 2, 3, 4, 5])
\end{lstlisting}

\subsection{Dictionary Keys}

\begin{lstlisting}
# Tuples can be dictionary keys (immutable)
coordinates_data = {
    (0, 0): "Origin",
    (1, 0): "Right",
    (0, 1): "Up",
    (1, 1): "Diagonal"
}

print(coordinates_data[(1, 1)])  # "Diagonal"
\end{lstlisting}

\section{Common Pitfalls}

\subsection{Single Element Tuple}

\begin{lstlisting}
# Wrong way (creates int, not tuple)
not_a_tuple = (42)
print(type(not_a_tuple))  # <class 'int'>

# Right way (comma is required)
single_tuple = (42,)
print(type(single_tuple))  # <class 'tuple'>
\end{lstlisting}

\subsection{Modifying Nested Mutable Objects}

\begin{lstlisting}
# Tuple is immutable, but contained objects might not be
tuple_with_list = ([1, 2, 3], "hello")

# This works (modifying the list inside the tuple)
tuple_with_list[0].append(4)
print(tuple_with_list)  # ([1, 2, 3, 4], 'hello')

# This doesn't work (trying to replace tuple element)
# tuple_with_list[0] = [5, 6, 7]  # TypeError
\end{lstlisting}

\section{Summary}

Tuples are essential Python data structures that provide:
\begin{itemize}
    \item \textbf{Immutability}: Cannot be changed after creation
    \item \textbf{Order}: Maintain insertion order
    \item \textbf{Hashability}: Can be used as dictionary keys
    \item \textbf{Memory efficiency}: More compact than lists
    \item \textbf{Multiple assignment}: Enable elegant unpacking
    \item \textbf{Data integrity}: Protect data from accidental modification
\end{itemize}

Use tuples when you need ordered, unchangeable collections of data!

\end{document}

\chapter{Dictionaries}
\documentclass[12pt,a4paper]{article}
\usepackage[utf8]{inputenc}
\usepackage{amsmath,amsfonts,amssymb}
\usepackage{listings}
\usepackage{xcolor}
\usepackage{geometry}
\usepackage{hyperref}
\usepackage{fancyhdr}

\geometry{margin=1in}
\pagestyle{fancy}
\fancyhf{}
\rhead{Python Dictionaries - Complete Guide}
\lfoot{Python for Data Science}
\rfoot{\thepage}

\lstset{
    language=Python,
    basicstyle=\ttfamily\small,
    keywordstyle=\color{blue}\bfseries,
    commentstyle=\color{green!60!black},
    stringstyle=\color{red},
    numbers=left,
    numberstyle=\tiny\color{gray},
    stepnumber=1,
    numbersep=8pt,
    showstringspaces=false,
    breaklines=true,
    frame=single,
    backgroundcolor=\color{gray!10}
}

\title{\textbf{Python Dictionaries - Complete Guide}}
\author{Python for Data Science - Beginner's Level}
\date{\today}

\begin{document}

\maketitle
\tableofcontents
\newpage

\section{Overview}

Dictionaries are \textbf{unordered} (Python 3.7+ maintains insertion order), \textbf{mutable}, and store data in \textbf{key-value pairs}. Keys must be \textbf{unique} and \textbf{immutable}, while values can be any data type.

\section{Dictionary Creation}

\subsection{Basic Creation}

\begin{lstlisting}
# Empty dictionary
empty_dict = {}
empty_dict = dict()

# Dictionary with initial values
student = {
    "name": "Alice",
    "age": 20,
    "grade": "A",
    "subjects": ["Math", "Physics", "Chemistry"]
}

# Mixed data types
mixed_dict = {
    1: "one",
    "two": 2,
    3.0: "three point zero",
    (4, 5): "tuple key"
}
\end{lstlisting}

\subsection{Using dict() Constructor}

\begin{lstlisting}
# From keyword arguments
person = dict(name="John", age=30, city="New York")

# From list of tuples
pairs = [("a", 1), ("b", 2), ("c", 3)]
from_pairs = dict(pairs)

# From zip
keys = ["name", "age", "city"]
values = ["Alice", 25, "Boston"]
from_zip = dict(zip(keys, values))
\end{lstlisting}

\subsection{Dictionary Comprehension}

\begin{lstlisting}
# Basic comprehension
squares = {x: x**2 for x in range(6)}

# With condition
even_squares = {x: x**2 for x in range(10) if x % 2 == 0}

# From existing dictionary
original = {"a": 1, "b": 2, "c": 3}
doubled = {k: v*2 for k, v in original.items()}
\end{lstlisting}

\section{Accessing Elements}

\subsection{Basic Access}

\begin{lstlisting}
student = {"name": "Alice", "age": 20, "grade": "A"}

# Direct access (raises KeyError if key doesn't exist)
print(student["name"])  # "Alice"
print(student["age"])   # 20

# Safe access with get()
print(student.get("name"))     # "Alice"
print(student.get("height"))   # None
print(student.get("height", "Unknown"))  # "Unknown"
\end{lstlisting}

\subsection{Advanced Access Methods}

\begin{lstlisting}
student = {"name": "Alice", "age": 20, "grade": "A"}

# Check if key exists
print("name" in student)        # True
print("height" in student)      # False

# Get all keys, values, items
print(student.keys())    # dict_keys(['name', 'age', 'grade'])
print(student.values())  # dict_values(['Alice', 20, 'A'])
print(student.items())   # dict_items([('name', 'Alice'), ...])
\end{lstlisting}

\section{Modifying Dictionaries}

\subsection{Adding and Updating Elements}

\begin{lstlisting}
student = {"name": "Alice", "age": 20}

# Add new key-value pair
student["grade"] = "A"
student["subjects"] = ["Math", "Physics"]

# Update existing value
student["age"] = 21

# Update multiple values
student.update({"grade": "A+", "city": "Boston"})

# Update from another dictionary
additional_info = {"phone": "123-456-7890", "email": "alice@email.com"}
student.update(additional_info)
\end{lstlisting}

\subsection{Removing Elements}

\begin{lstlisting}
student = {
    "name": "Alice", 
    "age": 21, 
    "grade": "A+", 
    "city": "Boston",
    "phone": "123-456-7890"
}

# pop() - Remove and return value
age = student.pop("age")

# pop() with default value
height = student.pop("height", "Not specified")

# popitem() - Remove and return last inserted item
last_item = student.popitem()

# del statement
del student["city"]

# clear() - Remove all elements
student.clear()
\end{lstlisting}

\section{Dictionary Methods}

\subsection{Comprehensive Method Overview}

\begin{lstlisting}
data = {"a": 1, "b": 2, "c": 3}

# copy() - Shallow copy
data_copy = data.copy()

# fromkeys() - Create dictionary with same value for all keys
keys = ["x", "y", "z"]
default_dict = dict.fromkeys(keys, 0)

# setdefault() - Get value or set default if key doesn't exist
data.setdefault("d", 4)  # Adds "d": 4
existing = data.setdefault("a", 100)  # Returns existing value
\end{lstlisting}

\section{Iterating Through Dictionaries}

\subsection{Basic Iteration}

\begin{lstlisting}
student = {"name": "Alice", "age": 21, "grade": "A+"}

# Iterate over keys (default behavior)
for key in student:
    print(f"{key}: {student[key]}")

# Explicitly iterate over keys
for key in student.keys():
    print(f"Key: {key}")

# Iterate over values
for value in student.values():
    print(f"Value: {value}")

# Iterate over key-value pairs
for key, value in student.items():
    print(f"{key}: {value}")
\end{lstlisting}

\section{Advanced Dictionary Operations}

\subsection{Nested Dictionaries}

\begin{lstlisting}
students = {
    "alice": {
        "age": 20,
        "grades": {"math": 95, "physics": 87, "chemistry": 92},
        "contact": {"email": "alice@email.com", "phone": "123-456-7890"}
    },
    "bob": {
        "age": 22,
        "grades": {"math": 78, "physics": 85, "chemistry": 80},
        "contact": {"email": "bob@email.com", "phone": "098-765-4321"}
    }
}

# Accessing nested values
print(students["alice"]["age"])  # 20
print(students["alice"]["grades"]["math"])  # 95
\end{lstlisting}

\subsection{Dictionary as Counter}

\begin{lstlisting}
# Manual counting
text = "hello world"
char_count = {}
for char in text:
    char_count[char] = char_count.get(char, 0) + 1

# Using collections.Counter (recommended)
from collections import Counter
char_counter = Counter(text)
\end{lstlisting}

\section{Dictionary Sorting and Filtering}

\subsection{Sorting Dictionaries}

\begin{lstlisting}
data = {"banana": 3, "apple": 5, "cherry": 1, "date": 4}

# Sort by keys
sorted_by_keys = dict(sorted(data.items()))

# Sort by values
sorted_by_values = dict(sorted(data.items(), key=lambda x: x[1]))

# Sort by values (descending)
sorted_desc = dict(sorted(data.items(), key=lambda x: x[1], reverse=True))
\end{lstlisting}

\subsection{Filtering Dictionaries}

\begin{lstlisting}
data = {"apple": 5, "banana": 3, "cherry": 8, "date": 2}

# Filter by value
high_values = {k: v for k, v in data.items() if v > 4}

# Filter by key
long_names = {k: v for k, v in data.items() if len(k) > 5}

# Filter using filter() function
even_values = dict(filter(lambda x: x[1] % 2 == 0, data.items()))
\end{lstlisting}

\section{Performance Considerations}

\begin{table}[h]
\centering
\begin{tabular}{|l|c|c|l|}
\hline
\textbf{Operation} & \textbf{Average} & \textbf{Worst} & \textbf{Notes} \\
\hline
Access/Search & O(1) & O(n) & Hash collision \\
Insert & O(1) & O(n) & Hash collision \\
Delete & O(1) & O(n) & Hash collision \\
Iteration & O(n) & O(n) & All elements \\
\hline
\end{tabular}
\caption{Dictionary Operations Time Complexity}
\end{table}

\section{Common Use Cases}

\subsection{Configuration Management}

\begin{lstlisting}
config = {
    "database": {
        "host": "localhost",
        "port": 5432,
        "name": "myapp"
    },
    "api": {
        "base_url": "https://api.example.com",
        "timeout": 30
    }
}
\end{lstlisting}

\subsection{Data Transformation}

\begin{lstlisting}
# Transform list of records to dictionary
students_list = [
    {"id": 1, "name": "Alice", "grade": "A"},
    {"id": 2, "name": "Bob", "grade": "B"}
]

# Index by ID
students_by_id = {student["id"]: student for student in students_list}
\end{lstlisting}

\section{Common Pitfalls}

\subsection{Mutable Default Arguments}

\begin{lstlisting}
# Wrong way
def add_item(item, target_dict={}):  # Dangerous!
    target_dict[len(target_dict)] = item
    return target_dict

# Right way
def add_item(item, target_dict=None):
    if target_dict is None:
        target_dict = {}
    target_dict[len(target_dict)] = item
    return target_dict
\end{lstlisting}

\subsection{Key Type Restrictions}

\begin{lstlisting}
# Valid keys (immutable types)
valid_dict = {
    "string": 1,
    42: 2,
    (1, 2): 3,
    frozenset([1, 2, 3]): 4
}

# Invalid keys (mutable types)
# invalid_dict = {
#     [1, 2, 3]: "list",      # TypeError
#     {"a": 1}: "dict",       # TypeError
# }
\end{lstlisting}

\section{Summary}

Dictionaries are powerful Python data structures that provide:
\begin{itemize}
    \item \textbf{Key-value mapping}: Efficient data organization
    \item \textbf{Fast lookup}: O(1) average case access time
    \item \textbf{Flexibility}: Support for various data types
    \item \textbf{Mutability}: Can be modified after creation
    \item \textbf{Rich methods}: Comprehensive built-in functionality
    \item \textbf{Memory efficiency}: Optimized hash table implementation
\end{itemize}

Master dictionaries to handle complex data relationships and build efficient applications!

\end{document}

\chapter{Sets}
\documentclass[12pt,a4paper]{article}
\usepackage[utf8]{inputenc}
\usepackage{amsmath,amsfonts,amssymb}
\usepackage{listings}
\usepackage{xcolor}
\usepackage{geometry}
\usepackage{hyperref}
\usepackage{fancyhdr}

\geometry{margin=1in}
\pagestyle{fancy}
\fancyhf{}
\rhead{Python Sets - Complete Guide}
\lfoot{Python for Data Science}
\rfoot{\thepage}

\lstset{
    language=Python,
    basicstyle=\ttfamily\small,
    keywordstyle=\color{blue}\bfseries,
    commentstyle=\color{green!60!black},
    stringstyle=\color{red},
    numbers=left,
    numberstyle=\tiny\color{gray},
    stepnumber=1,
    numbersep=8pt,
    showstringspaces=false,
    breaklines=true,
    frame=single,
    backgroundcolor=\color{gray!10}
}

\title{\textbf{Python Sets - Complete Guide}}
\author{Python for Data Science - Beginner's Level}
\date{\today}

\begin{document}

\maketitle
\tableofcontents
\newpage

\section{Overview}

Sets are \textbf{unordered}, \textbf{mutable} collections of \textbf{unique} elements. They are ideal for membership testing, removing duplicates, and mathematical set operations like union, intersection, and difference.

\section{Set Creation}

\subsection{Basic Creation}

\begin{lstlisting}
# Empty set (note: {} creates an empty dict, not set)
empty_set = set()

# Set with initial values
numbers = {1, 2, 3, 4, 5}
mixed_set = {1, "hello", 3.14, True}

# Note: duplicates are automatically removed
duplicate_values = {1, 2, 2, 3, 3, 4}
print(duplicate_values)  # {1, 2, 3, 4}

# Set from string (each character becomes an element)
char_set = set("hello")
print(char_set)  # {'h', 'e', 'l', 'o'}
\end{lstlisting}

\subsection{Using set() Constructor}

\begin{lstlisting}
# From list
from_list = set([1, 2, 3, 2, 1])  # {1, 2, 3}

# From tuple
from_tuple = set((1, 2, 3, 2, 1))  # {1, 2, 3}

# From string
from_string = set("programming")  # {'p', 'r', 'o', 'g', 'a', 'm', 'i', 'n'}

# From range
from_range = set(range(5))  # {0, 1, 2, 3, 4}
\end{lstlisting}

\subsection{Set Comprehension}

\begin{lstlisting}
# Basic comprehension
squares = {x**2 for x in range(6)}  # {0, 1, 4, 9, 16, 25}

# With condition
even_squares = {x**2 for x in range(10) if x % 2 == 0}

# From existing iterable
words = ["apple", "banana", "apple", "cherry"]
unique_lengths = {len(word) for word in words}  # {5, 6}
\end{lstlisting}

\section{Adding and Removing Elements}

\subsection{Adding Elements}

\begin{lstlisting}
fruits = {"apple", "banana"}

# add() - Add single element
fruits.add("cherry")

# Adding existing element (no effect)
fruits.add("apple")  # No change

# update() - Add multiple elements
fruits.update(["orange", "grape"])

# Using |= operator (union update)
more_fruits = {"pineapple", "strawberry"}
fruits |= more_fruits
\end{lstlisting}

\subsection{Removing Elements}

\begin{lstlisting}
fruits = {"apple", "banana", "cherry", "orange", "grape"}

# remove() - Remove element (raises KeyError if not found)
fruits.remove("banana")

# discard() - Remove element (no error if not found)
fruits.discard("cherry")
fruits.discard("mango")  # No error

# pop() - Remove and return arbitrary element
removed = fruits.pop()

# clear() - Remove all elements
fruits.clear()
\end{lstlisting}

\section{Set Operations and Methods}

\subsection{Membership Testing}

\begin{lstlisting}
numbers = {1, 2, 3, 4, 5}

# in operator (very fast - O(1) average case)
print(3 in numbers)      # True
print(6 in numbers)      # False
print(3 not in numbers)  # False
\end{lstlisting}

\subsection{Mathematical Set Operations}

\subsubsection{Union Operations}

\begin{lstlisting}
set1 = {1, 2, 3, 4}
set2 = {3, 4, 5, 6}

# union() method
union_result = set1.union(set2)  # {1, 2, 3, 4, 5, 6}

# | operator
union_operator = set1 | set2  # {1, 2, 3, 4, 5, 6}

# In-place union
set1_copy = set1.copy()
set1_copy.update(set2)  # or set1_copy |= set2
\end{lstlisting}

\subsubsection{Intersection Operations}

\begin{lstlisting}
set1 = {1, 2, 3, 4, 5}
set2 = {3, 4, 5, 6, 7}

# intersection() method
intersection_result = set1.intersection(set2)  # {3, 4, 5}

# & operator
intersection_operator = set1 & set2  # {3, 4, 5}

# In-place intersection
set1_copy = set1.copy()
set1_copy.intersection_update(set2)  # or set1_copy &= set2
\end{lstlisting}

\subsubsection{Difference Operations}

\begin{lstlisting}
set1 = {1, 2, 3, 4, 5}
set2 = {3, 4, 5, 6, 7}

# difference() method (elements in set1 but not in set2)
difference_result = set1.difference(set2)  # {1, 2}

# - operator
difference_operator = set1 - set2  # {1, 2}

# Reverse difference
reverse_diff = set2 - set1  # {6, 7}
\end{lstlisting}

\subsubsection{Symmetric Difference Operations}

\begin{lstlisting}
set1 = {1, 2, 3, 4}
set2 = {3, 4, 5, 6}

# symmetric_difference() method (elements in either set, but not both)
sym_diff_result = set1.symmetric_difference(set2)  # {1, 2, 5, 6}

# ^ operator
sym_diff_operator = set1 ^ set2  # {1, 2, 5, 6}
\end{lstlisting}

\section{Set Relationship Testing}

\begin{lstlisting}
set1 = {1, 2, 3}
set2 = {1, 2, 3, 4, 5}
set3 = {4, 5, 6}

# issubset() - Check if set is subset of another
print(set1.issubset(set2))  # True
print(set1 <= set2)         # True (alternative syntax)

# issuperset() - Check if set is superset of another
print(set2.issuperset(set1))  # True
print(set2 >= set1)           # True (alternative syntax)

# isdisjoint() - Check if sets have no common elements
print(set1.isdisjoint(set3))  # True
\end{lstlisting}

\section{Advanced Set Operations}

\subsection{Frozen Sets}

\begin{lstlisting}
# Immutable sets - can be used as dictionary keys
regular_set = {1, 2, 3}
frozen_set = frozenset([1, 2, 3])

# Frozen sets can be dictionary keys
dict_with_frozenset_keys = {
    frozenset([1, 2]): "first",
    frozenset([3, 4]): "second"
}

# Frozen sets support all non-mutating operations
fs1 = frozenset([1, 2, 3])
fs2 = frozenset([2, 3, 4])

print(fs1 | fs2)  # frozenset({1, 2, 3, 4})
print(fs1 & fs2)  # frozenset({2, 3})
\end{lstlisting}

\section{Performance Considerations}

\begin{table}[h]
\centering
\begin{tabular}{|l|c|c|l|}
\hline
\textbf{Operation} & \textbf{Average} & \textbf{Worst} & \textbf{Notes} \\
\hline
Add & O(1) & O(n) & Hash collision \\
Remove & O(1) & O(n) & Hash collision \\
Search (in) & O(1) & O(n) & Hash collision \\
Union & O(len(s1) + len(s2)) & O(len(s1) * len(s2)) & \\
Intersection & O(min(len(s1), len(s2))) & O(len(s1) * len(s2)) & \\
\hline
\end{tabular}
\caption{Set Operations Time Complexity}
\end{table}

\section{Set Applications and Use Cases}

\subsection{Removing Duplicates}

\begin{lstlisting}
# Remove duplicates from list
numbers = [1, 2, 2, 3, 1, 4, 3, 5]
unique_numbers = list(set(numbers))

# Preserve order while removing duplicates
def remove_duplicates_ordered(lst):
    seen = set()
    result = []
    for item in lst:
        if item not in seen:
            seen.add(item)
            result.append(item)
    return result
\end{lstlisting}

\subsection{Finding Common Elements}

\begin{lstlisting}
# Find common elements across multiple lists
def find_common_elements(*lists):
    if not lists:
        return set()
    
    common = set(lists[0])
    for lst in lists[1:]:
        common &= set(lst)
    return common

# Test
list1 = [1, 2, 3, 4, 5]
list2 = [3, 4, 5, 6, 7]
list3 = [4, 5, 6, 7, 8]

common = find_common_elements(list1, list2, list3)  # {4, 5}
\end{lstlisting}

\subsection{Data Validation}

\begin{lstlisting}
# Validate data against allowed values
ALLOWED_COLORS = {"red", "green", "blue", "yellow", "orange"}
ALLOWED_SIZES = {"small", "medium", "large", "xl", "xxl"}

def validate_product(color, size):
    errors = []
    
    if color.lower() not in ALLOWED_COLORS:
        errors.append(f"Invalid color: {color}")
    
    if size.lower() not in ALLOWED_SIZES:
        errors.append(f"Invalid size: {size}")
    
    return errors
\end{lstlisting}

\section{Common Pitfalls}

\subsection{Mutable Elements in Sets}

\begin{lstlisting}
# This won't work - lists are mutable and unhashable
# invalid_set = {[1, 2], [3, 4]}  # TypeError

# Use tuples or frozensets instead
valid_set = {(1, 2), (3, 4)}
frozenset_set = {frozenset([1, 2]), frozenset([3, 4])}
\end{lstlisting}

\subsection{Empty Set Creation}

\begin{lstlisting}
# Wrong way to create empty set
empty_dict = {}  # This creates a dictionary, not a set!

# Right way to create empty set
empty_set = set()
\end{lstlisting}

\section{Summary}

Sets are powerful Python data structures that provide:
\begin{itemize}
    \item \textbf{Uniqueness}: Automatically handle duplicate elimination
    \item \textbf{Fast membership testing}: O(1) average case lookup
    \item \textbf{Mathematical operations}: Union, intersection, difference, etc.
    \item \textbf{Unordered collections}: Focus on membership, not position
    \item \textbf{Hashable elements only}: Immutable objects as elements
    \item \textbf{Memory efficiency}: Optimized for unique element storage
\end{itemize}

Master sets to efficiently handle unique collections and perform mathematical set operations!

\end{document}

\chapter{NumPy Arrays}
\documentclass[12pt,a4paper]{article}
\usepackage[utf8]{inputenc}
\usepackage{amsmath,amsfonts,amssymb}
\usepackage{listings}
\usepackage{xcolor}
\usepackage{geometry}
\usepackage{hyperref}
\usepackage{fancyhdr}

\geometry{margin=1in}
\pagestyle{fancy}
\fancyhf{}
\rhead{NumPy - Complete Guide}
\lfoot{Python for Data Science}
\rfoot{\thepage}

\lstset{
    language=Python,
    basicstyle=\ttfamily\small,
    keywordstyle=\color{blue}\bfseries,
    commentstyle=\color{green!60!black},
    stringstyle=\color{red},
    numbers=left,
    numberstyle=\tiny\color{gray},
    stepnumber=1,
    numbersep=8pt,
    showstringspaces=false,
    breaklines=true,
    frame=single,
    backgroundcolor=\color{gray!10}
}

\title{\textbf{NumPy - Complete Guide}}
\author{Python for Data Science - Beginner's Level}
\date{\today}

\begin{document}

\maketitle
\tableofcontents
\newpage

\section{Overview}

NumPy (Numerical Python) is the fundamental package for scientific computing in Python. It provides powerful N-dimensional array objects and tools for working with arrays.

\section{Installation and Import}

\begin{lstlisting}
# Installation
# pip install numpy

# Import
import numpy as np
\end{lstlisting}

\section{Array Creation}

\subsection{Basic Array Creation}

\begin{lstlisting}
# From lists
arr1d = np.array([1, 2, 3, 4, 5])
arr2d = np.array([[1, 2, 3], [4, 5, 6]])
arr3d = np.array([[[1, 2], [3, 4]], [[5, 6], [7, 8]]])

# Array properties
print(arr2d.shape)    # (2, 3)
print(arr2d.dtype)    # int64
print(arr2d.ndim)     # 2
print(arr2d.size)     # 6
\end{lstlisting}

\subsection{Specialized Array Creation}

\begin{lstlisting}
# Zeros and ones
zeros = np.zeros((3, 4))           # 3x4 array of zeros
ones = np.ones((2, 3))             # 2x3 array of ones
full = np.full((2, 2), 7)          # 2x2 array filled with 7

# Identity and diagonal
identity = np.eye(3)               # 3x3 identity matrix
diag = np.diag([1, 2, 3])         # Diagonal matrix

# Range arrays
arange = np.arange(0, 10, 2)       # [0, 2, 4, 6, 8]
linspace = np.linspace(0, 1, 5)    # [0, 0.25, 0.5, 0.75, 1]

# Random arrays
random = np.random.rand(3, 3)      # Random values [0, 1)
randint = np.random.randint(1, 10, (2, 3))  # Random integers
\end{lstlisting}

\section{Array Indexing and Slicing}

\subsection{Basic Indexing}

\begin{lstlisting}
arr = np.array([[1, 2, 3, 4],
                [5, 6, 7, 8],
                [9, 10, 11, 12]])

# Single element
print(arr[1, 2])      # 7

# Row and column
print(arr[1, :])      # [5 6 7 8] (row 1)
print(arr[:, 2])      # [3 7 11] (column 2)

# Slicing
print(arr[0:2, 1:3])  # [[2 3], [6 7]]
\end{lstlisting}

\subsection{Advanced Indexing}

\begin{lstlisting}
arr = np.array([1, 2, 3, 4, 5, 6, 7, 8, 9, 10])

# Boolean indexing
mask = arr > 5
print(arr[mask])      # [6 7 8 9 10]

# Fancy indexing
indices = [1, 3, 5]
print(arr[indices])   # [2 4 6]

# Conditional selection
result = np.where(arr > 5, arr, 0)  # Replace values ≤5 with 0
\end{lstlisting}

\section{Array Operations}

\subsection{Mathematical Operations}

\begin{lstlisting}
arr1 = np.array([1, 2, 3, 4])
arr2 = np.array([5, 6, 7, 8])

# Element-wise operations
print(arr1 + arr2)    # [6 8 10 12]
print(arr1 * arr2)    # [5 12 21 32]
print(arr1 ** 2)      # [1 4 9 16]

# Universal functions
print(np.sqrt(arr1))  # [1. 1.414 1.732 2.]
print(np.exp(arr1))   # [2.718 7.389 20.086 54.598]
print(np.sin(arr1))   # [0.841 0.909 0.141 -0.757]
\end{lstlisting}

\subsection{Aggregation Functions}

\begin{lstlisting}
arr = np.array([[1, 2, 3],
                [4, 5, 6]])

# Basic aggregations
print(np.sum(arr))        # 21
print(np.mean(arr))       # 3.5
print(np.std(arr))        # 1.707
print(np.min(arr))        # 1
print(np.max(arr))        # 6

# Axis-specific aggregations
print(np.sum(arr, axis=0))    # [5 7 9] (column sums)
print(np.sum(arr, axis=1))    # [6 15] (row sums)
\end{lstlisting}

\section{Array Manipulation}

\subsection{Reshaping}

\begin{lstlisting}
arr = np.arange(12)  # [0 1 2 3 4 5 6 7 8 9 10 11]

# Reshape
reshaped = arr.reshape(3, 4)

# Flatten
flattened = reshaped.flatten()

# Transpose
transposed = reshaped.T
\end{lstlisting}

\subsection{Joining and Splitting}

\begin{lstlisting}
arr1 = np.array([[1, 2], [3, 4]])
arr2 = np.array([[5, 6], [7, 8]])

# Concatenation
hstack = np.hstack([arr1, arr2])  # [[1 2 5 6], [3 4 7 8]]
vstack = np.vstack([arr1, arr2])  # [[1 2], [3 4], [5 6], [7 8]]

# Splitting
arr = np.arange(8)
split = np.split(arr, 4)  # [array([0, 1]), array([2, 3]), ...]
\end{lstlisting}

\section{Linear Algebra}

\begin{lstlisting}
# Matrix operations
A = np.array([[1, 2], [3, 4]])
B = np.array([[5, 6], [7, 8]])

# Matrix multiplication
dot_product = np.dot(A, B)    # or A @ B

# Matrix properties
det = np.linalg.det(A)        # Determinant
inv = np.linalg.inv(A)        # Inverse
eigenvals = np.linalg.eigvals(A)  # Eigenvalues

# Solving linear systems Ax = b
b = np.array([1, 2])
x = np.linalg.solve(A, b)
\end{lstlisting}

\section{Broadcasting}

\begin{lstlisting}
# Broadcasting allows operations between arrays of different shapes
arr = np.array([[1, 2, 3],
                [4, 5, 6]])

# Scalar broadcasting
result1 = arr + 10            # Add 10 to all elements

# Vector broadcasting
vec = np.array([10, 20, 30])
result2 = arr + vec           # Add vector to each row

# Column vector broadcasting
col_vec = np.array([[10], [20]])
result3 = arr + col_vec       # Add column vector
\end{lstlisting}

\section{Performance Tips}

\begin{lstlisting}
# Use vectorized operations instead of loops
# Slow
result_slow = []
for i in range(len(arr)):
    result_slow.append(arr[i] ** 2)

# Fast
result_fast = arr ** 2

# Use appropriate data types
int_arr = np.array([1, 2, 3], dtype=np.int32)    # Less memory
float_arr = np.array([1, 2, 3], dtype=np.float64) # More precision
\end{lstlisting}

\section{Summary}

NumPy provides:
\begin{itemize}
    \item \textbf{Efficient arrays}: Fast numerical operations
    \item \textbf{Broadcasting}: Operations between different shapes
    \item \textbf{Linear algebra}: Matrix operations and decompositions
    \item \textbf{Random numbers}: Statistical distributions and sampling
    \item \textbf{Universal functions}: Element-wise operations
    \item \textbf{Memory efficiency}: Optimized C implementations
\end{itemize}

Essential for scientific computing and data science in Python!

\end{document}

\chapter{Pandas DataFrames}
\documentclass[12pt,a4paper]{article}
\usepackage[utf8]{inputenc}
\usepackage{amsmath,amsfonts,amssymb}
\usepackage{listings}
\usepackage{xcolor}
\usepackage{geometry}
\usepackage{hyperref}
\usepackage{fancyhdr}

\geometry{margin=1in}
\pagestyle{fancy}
\fancyhf{}
\rhead{Pandas - Complete Guide}
\lfoot{Python for Data Science}
\rfoot{\thepage}

\lstset{
    language=Python,
    basicstyle=\ttfamily\small,
    keywordstyle=\color{blue}\bfseries,
    commentstyle=\color{green!60!black},
    stringstyle=\color{red},
    numbers=left,
    numberstyle=\tiny\color{gray},
    stepnumber=1,
    numbersep=8pt,
    showstringspaces=false,
    breaklines=true,
    frame=single,
    backgroundcolor=\color{gray!10}
}

\title{\textbf{Pandas - Complete Guide}}
\author{Python for Data Science - Beginner's Level}
\date{\today}

\begin{document}

\maketitle
\tableofcontents
\newpage

\section{Overview}

Pandas is a powerful data manipulation and analysis library built on NumPy. It provides DataFrame and Series objects for handling structured data.

\section{Installation and Import}

\begin{lstlisting}
# Installation
# pip install pandas

# Import
import pandas as pd
import numpy as np
\end{lstlisting}

\section{Core Data Structures}

\subsection{Series}

\begin{lstlisting}
# Create Series
s1 = pd.Series([1, 2, 3, 4, 5])
s2 = pd.Series([1, 2, 3], index=['a', 'b', 'c'])
s3 = pd.Series({'A': 1, 'B': 2, 'C': 3})

# Series properties
print(s2.values)      # [1 2 3]
print(s2.index)       # Index(['a', 'b', 'c'])
print(s2.dtype)       # int64
print(s2.shape)       # (3,)
\end{lstlisting}

\subsection{DataFrame}

\begin{lstlisting}
# Create DataFrame
data = {
    'Name': ['Alice', 'Bob', 'Charlie', 'Diana'],
    'Age': [25, 30, 35, 28],
    'City': ['NYC', 'LA', 'Chicago', 'Boston'],
    'Salary': [50000, 60000, 70000, 55000]
}
df = pd.DataFrame(data)

# DataFrame properties
print(df.shape)       # (4, 4)
print(df.columns)     # Index(['Name', 'Age', 'City', 'Salary'])
print(df.index)       # RangeIndex(start=0, stop=4, step=1)
print(df.dtypes)      # Data types of each column
\end{lstlisting}

\section{Data Selection and Indexing}

\subsection{Basic Selection}

\begin{lstlisting}
# Column selection
names = df['Name']              # Series
subset = df[['Name', 'Age']]    # DataFrame

# Row selection
first_row = df.iloc[0]          # First row by position
named_row = df.loc[0]           # First row by label

# Slicing
first_three = df.iloc[0:3]      # First 3 rows
columns_slice = df.iloc[:, 1:3] # All rows, columns 1-2
\end{lstlisting}

\subsection{Advanced Selection}

\begin{lstlisting}
# Boolean indexing
high_salary = df[df['Salary'] > 55000]
young_high_earners = df[(df['Age'] < 30) & (df['Salary'] > 50000)]

# Query method
filtered = df.query('Age > 25 and Salary < 65000')

# isin method
cities_filter = df[df['City'].isin(['NYC', 'LA'])]

# String methods
name_filter = df[df['Name'].str.startswith('A')]
\end{lstlisting}

\section{Data Manipulation}

\subsection{Adding and Modifying Data}

\begin{lstlisting}
# Add new column
df['Bonus'] = df['Salary'] * 0.1
df['Full_Info'] = df['Name'] + ' - ' + df['City']

# Modify existing data
df.loc[df['Age'] > 30, 'Category'] = 'Senior'
df.loc[df['Age'] <= 30, 'Category'] = 'Junior'

# Apply functions
df['Age_Group'] = df['Age'].apply(lambda x: 'Young' if x < 30 else 'Mature')
\end{lstlisting}

\subsection{Removing Data}

\begin{lstlisting}
# Drop columns
df_no_bonus = df.drop('Bonus', axis=1)
df_subset = df.drop(['Bonus', 'Full_Info'], axis=1)

# Drop rows
df_no_first = df.drop(0, axis=0)
df_filtered = df.drop(df[df['Age'] < 25].index)

# Remove duplicates
df_unique = df.drop_duplicates()
df_unique_names = df.drop_duplicates(subset=['Name'])
\end{lstlisting}

\section{Data Cleaning}

\subsection{Handling Missing Data}

\begin{lstlisting}
# Create data with missing values
df_missing = df.copy()
df_missing.loc[1, 'Salary'] = np.nan
df_missing.loc[2, 'Age'] = np.nan

# Check for missing data
print(df_missing.isnull().sum())
print(df_missing.info())

# Handle missing data
df_dropped = df_missing.dropna()                    # Drop rows with any NaN
df_filled = df_missing.fillna(0)                    # Fill with 0
df_mean = df_missing.fillna(df_missing.mean())      # Fill with mean
\end{lstlisting}

\subsection{Data Type Conversion}

\begin{lstlisting}
# Convert data types
df['Age'] = df['Age'].astype('int32')
df['Salary'] = df['Salary'].astype('float64')

# Convert to datetime
dates = pd.Series(['2023-01-01', '2023-02-01', '2023-03-01'])
df['Date'] = pd.to_datetime(dates)

# Convert to categorical
df['City'] = df['City'].astype('category')
\end{lstlisting}

\section{Data Analysis}

\subsection{Descriptive Statistics}

\begin{lstlisting}
# Basic statistics
print(df.describe())                    # Summary statistics
print(df['Salary'].mean())              # Mean salary
print(df['Age'].median())               # Median age
print(df['City'].value_counts())        # Count by city

# Correlation
numeric_df = df.select_dtypes(include=[np.number])
correlation_matrix = numeric_df.corr()
\end{lstlisting}

\subsection{Grouping and Aggregation}

\begin{lstlisting}
# Group by single column
city_groups = df.groupby('City')
print(city_groups['Salary'].mean())     # Average salary by city
print(city_groups.size())               # Count by city

# Group by multiple columns
df['Age_Group'] = pd.cut(df['Age'], bins=[0, 30, 40, 100], 
                        labels=['Young', 'Middle', 'Senior'])
multi_group = df.groupby(['City', 'Age_Group'])
print(multi_group['Salary'].agg(['mean', 'count']))

# Custom aggregations
agg_result = df.groupby('City').agg({
    'Salary': ['mean', 'max', 'min'],
    'Age': 'mean'
})
\end{lstlisting}

\subsection{Pivot Tables}

\begin{lstlisting}
# Create pivot table
pivot = df.pivot_table(
    values='Salary',
    index='City',
    columns='Age_Group',
    aggfunc='mean',
    fill_value=0
)

# Cross-tabulation
crosstab = pd.crosstab(df['City'], df['Age_Group'])
\end{lstlisting}

\section{Data Merging and Joining}

\subsection{Concatenation}

\begin{lstlisting}
# Create additional data
df2 = pd.DataFrame({
    'Name': ['Eve', 'Frank'],
    'Age': [32, 29],
    'City': ['Seattle', 'Miami'],
    'Salary': [65000, 58000]
})

# Concatenate vertically
combined = pd.concat([df, df2], ignore_index=True)

# Concatenate horizontally
df_extra = pd.DataFrame({'Department': ['IT', 'HR', 'Finance', 'Marketing']})
horizontal = pd.concat([df, df_extra], axis=1)
\end{lstlisting}

\subsection{Merging}

\begin{lstlisting}
# Create related data
departments = pd.DataFrame({
    'Name': ['Alice', 'Bob', 'Charlie', 'Diana'],
    'Department': ['IT', 'HR', 'Finance', 'Marketing']
})

# Inner join
merged_inner = pd.merge(df, departments, on='Name', how='inner')

# Left join
merged_left = pd.merge(df, departments, on='Name', how='left')
\end{lstlisting}

\section{File I/O Operations}

\subsection{Reading Data}

\begin{lstlisting}
# CSV files
df_csv = pd.read_csv('data.csv')
df_csv_custom = pd.read_csv('data.csv', sep=';', header=0, index_col=0)

# Excel files
df_excel = pd.read_excel('data.xlsx', sheet_name='Sheet1')

# JSON files
df_json = pd.read_json('data.json')
\end{lstlisting}

\subsection{Writing Data}

\begin{lstlisting}
# CSV files
df.to_csv('output.csv', index=False)
df.to_csv('output_custom.csv', sep=';', header=True)

# Excel files
df.to_excel('output.xlsx', sheet_name='Data', index=False)

# JSON files
df.to_json('output.json', orient='records')
\end{lstlisting}

\section{Time Series Analysis}

\begin{lstlisting}
# Create time series data
dates = pd.date_range('2023-01-01', periods=100, freq='D')
ts_data = pd.DataFrame({
    'Date': dates,
    'Value': np.random.randn(100).cumsum()
})
ts_data.set_index('Date', inplace=True)

# Time-based selection
jan_data = ts_data['2023-01']
week_data = ts_data['2023-01-01':'2023-01-07']

# Resampling
monthly = ts_data.resample('M').mean()      # Monthly average
weekly = ts_data.resample('W').sum()        # Weekly sum

# Rolling operations
ts_data['Rolling_Mean'] = ts_data['Value'].rolling(window=7).mean()
\end{lstlisting}

\section{Performance Tips}

\begin{lstlisting}
# Use vectorized operations
df['Total'] = df['Quantity'] * df['Price']  # Fast

# Use appropriate data types
df['Category'] = df['Category'].astype('category')  # Memory efficient

# Use query for complex filtering
result = df.query('Price > 100 and Category == "Electronics"')  # Fast

# Use chaining for multiple operations
result = (df
          .query('Age > 25')
          .groupby('City')['Salary']
          .mean()
          .sort_values(ascending=False))
\end{lstlisting}

\section{Summary}

Pandas provides:
\begin{itemize}
    \item \textbf{DataFrame/Series}: Powerful data structures
    \item \textbf{Data cleaning}: Handle missing values and duplicates
    \item \textbf{Data analysis}: Grouping, aggregation, and statistics
    \item \textbf{File I/O}: Read/write various formats
    \item \textbf{Time series}: Date/time operations and resampling
    \item \textbf{Merging}: Join and combine datasets
    \item \textbf{Performance}: Optimized operations for large datasets
\end{itemize}

Essential for data analysis and manipulation in Python!

\end{document}

\chapter{Best Practices and Guidelines}

\section{Choosing the Right Data Structure}

\subsection{When to Use Lists}
\begin{itemize}
    \item Need ordered, mutable collections
    \item Frequent appending of elements
    \item Need to maintain insertion order
    \item Working with sequences of related items
\end{itemize}

\subsection{When to Use Tuples}
\begin{itemize}
    \item Need immutable sequences
    \item Returning multiple values from functions
    \item Using as dictionary keys
    \item Representing fixed collections (coordinates, RGB values)
\end{itemize}

\subsection{When to Use Dictionaries}
\begin{itemize}
    \item Need key-value associations
    \item Fast lookups by key
    \item Representing structured data
    \item Caching and memoization
\end{itemize}

\subsection{When to Use Sets}
\begin{itemize}
    \item Need unique collections
    \item Fast membership testing
    \item Mathematical set operations
    \item Removing duplicates
\end{itemize}

\subsection{When to Use NumPy}
\begin{itemize}
    \item Numerical computations
    \item Large arrays of homogeneous data
    \item Mathematical operations on arrays
    \item Scientific computing
\end{itemize}

\subsection{When to Use Pandas}
\begin{itemize}
    \item Data analysis and manipulation
    \item Working with structured/tabular data
    \item Data cleaning and preprocessing
    \item Time series analysis
\end{itemize}

\section{Performance Optimization Tips}

\begin{enumerate}
    \item \textbf{Use appropriate data types}: Choose the most suitable data structure for your use case
    \item \textbf{Leverage built-in functions}: Use built-in methods and functions when possible
    \item \textbf{Avoid premature optimization}: Profile your code before optimizing
    \item \textbf{Use list comprehensions}: More efficient than traditional loops for simple operations
    \item \textbf{Consider memory usage}: Be mindful of memory consumption with large datasets
    \item \textbf{Use vectorized operations}: Prefer NumPy/Pandas vectorized operations over loops
\end{enumerate}

\section{Common Pitfalls to Avoid}

\begin{enumerate}
    \item \textbf{Mutable default arguments}: Don't use mutable objects as default function arguments
    \item \textbf{Modifying collections during iteration}: Create copies or iterate backwards when modifying
    \item \textbf{Inefficient string concatenation}: Use join() for multiple string concatenations
    \item \textbf{Not handling missing data}: Always check for and handle missing/null values
    \item \textbf{Ignoring data types}: Be explicit about data types, especially in NumPy/Pandas
\end{enumerate}

\chapter{Conclusion}

Mastering Python data structures is essential for effective programming and data science. Each structure has its strengths and optimal use cases:

\begin{itemize}
    \item \textbf{Lists} provide flexibility and mutability for dynamic collections
    \item \textbf{Tuples} offer immutability and memory efficiency for fixed data
    \item \textbf{Dictionaries} enable fast key-based lookups and data organization
    \item \textbf{Sets} handle unique collections and mathematical operations efficiently
    \item \textbf{NumPy} powers scientific computing with efficient array operations
    \item \textbf{Pandas} streamlines data analysis and manipulation workflows
\end{itemize}

Continue practicing with these data structures, experiment with different approaches, and always consider the specific requirements of your problem when choosing the appropriate structure.

\textbf{Next Steps:}
\begin{itemize}
    \item Practice with real datasets
    \item Explore advanced features of each library
    \item Learn about specialized data structures in the collections module
    \item Study algorithm complexity and optimization techniques
    \item Apply these concepts to data science projects
\end{itemize}

\appendix

\chapter{Quick Reference}

\section{List Methods}
\begin{lstlisting}
# Common list methods
list.append(item)       # Add item to end
list.insert(i, item)    # Insert item at position i
list.remove(item)       # Remove first occurrence
list.pop(i)            # Remove and return item at position i
list.index(item)       # Find index of item
list.count(item)       # Count occurrences
list.sort()            # Sort in place
list.reverse()         # Reverse in place
\end{lstlisting}

\section{Dictionary Methods}
\begin{lstlisting}
# Common dictionary methods
dict.get(key, default)     # Safe key access
dict.keys()               # Get all keys
dict.values()             # Get all values
dict.items()              # Get key-value pairs
dict.pop(key, default)    # Remove and return value
dict.update(other)        # Update with another dict
dict.setdefault(key, val) # Set default if key missing
\end{lstlisting}

\section{Set Operations}
\begin{lstlisting}
# Set operations
set1.union(set2)              # set1 | set2
set1.intersection(set2)       # set1 & set2
set1.difference(set2)         # set1 - set2
set1.symmetric_difference(set2) # set1 ^ set2
set1.issubset(set2)          # set1 <= set2
set1.issuperset(set2)        # set1 >= set2
\end{lstlisting}

\section{NumPy Functions}
\begin{lstlisting}
# Essential NumPy functions
np.array(data)           # Create array
np.zeros(shape)          # Array of zeros
np.ones(shape)           # Array of ones
np.arange(start, stop)   # Range array
np.linspace(start, stop, num) # Linear spacing
np.reshape(arr, shape)   # Reshape array
np.sum(arr, axis)        # Sum along axis
np.mean(arr, axis)       # Mean along axis
\end{lstlisting}

\section{Pandas Operations}
\begin{lstlisting}
# Essential Pandas operations
pd.DataFrame(data)       # Create DataFrame
df.head(n)              # First n rows
df.tail(n)              # Last n rows
df.info()               # DataFrame info
df.describe()           # Summary statistics
df.groupby(column)      # Group by column
df.merge(other)         # Merge DataFrames
df.to_csv(filename)     # Export to CSV
\end{lstlisting}

\end{document}